\documentclass[a4paper,12pt,twoside]{memoir}

% Castellano
\usepackage[spanish,es-tabla]{babel}
\selectlanguage{spanish}
\usepackage[utf8]{inputenc}
\usepackage[T1]{fontenc}
\usepackage{lmodern} % Scalable font
\usepackage{microtype}
\usepackage{placeins}

\RequirePackage{booktabs}
\RequirePackage[table]{xcolor}
\RequirePackage{xtab}
\RequirePackage{multirow}

% Links
\PassOptionsToPackage{hyphens}{url}\usepackage[colorlinks]{hyperref}
\hypersetup{
	allcolors = {red}
}

% Ecuaciones
\usepackage{amsmath}

% Rutas de fichero / paquete
\newcommand{\ruta}[1]{{\sffamily #1}}

% Párrafos
\nonzeroparskip

% Huérfanas y viudas
\widowpenalty100000
\clubpenalty100000

% Imagenes
\usepackage{graphicx}
\newcommand{\imagen}[2]{
	\begin{figure}[!h]
		\centering
		\includegraphics[width=0.9\textwidth]{#1}
		\caption{#2}\label{fig:#1}
	\end{figure}
	\FloatBarrier
}

\newcommand{\imagenflotante}[2]{
	\begin{figure}%[!h]
		\centering
		\includegraphics[width=0.9\textwidth]{#1}
		\caption{#2}\label{fig:#1}
	\end{figure}
}



% El comando \figura nos permite insertar figuras comodamente, y utilizando
% siempre el mismo formato. Los parametros son:
% 1 -> Porcentaje del ancho de página que ocupará la figura (de 0 a 1)
% 2 --> Fichero de la imagen
% 3 --> Texto a pie de imagen
% 4 --> Etiqueta (label) para referencias
% 5 --> Opciones que queramos pasarle al \includegraphics
% 6 --> Opciones de posicionamiento a pasarle a \begin{figure}
\newcommand{\figuraConPosicion}[6]{%
  \setlength{\anchoFloat}{#1\textwidth}%
  \addtolength{\anchoFloat}{-4\fboxsep}%
  \setlength{\anchoFigura}{\anchoFloat}%
  \begin{figure}[#6]
    \begin{center}%
      \Ovalbox{%
        \begin{minipage}{\anchoFloat}%
          \begin{center}%
            \includegraphics[width=\anchoFigura,#5]{#2}%
            \caption{#3}%
            \label{#4}%
          \end{center}%
        \end{minipage}
      }%
    \end{center}%
  \end{figure}%
}

%
% Comando para incluir imágenes en formato apaisado (sin marco).
\newcommand{\figuraApaisadaSinMarco}[5]{%
  \begin{figure}%
    \begin{center}%
    \includegraphics[angle=90,height=#1\textheight,#5]{#2}%
    \caption{#3}%
    \label{#4}%
    \end{center}%
  \end{figure}%
}
% Para las tablas
\newcommand{\otoprule}{\midrule [\heavyrulewidth]}
%
% Nuevo comando para tablas pequeñas (menos de una página).
\newcommand{\tablaSmall}[5]{%
 \begin{table}
  \begin{center}
   \rowcolors {2}{gray!35}{}
   \begin{tabular}{#2}
    \toprule
    #4
    \otoprule
    #5
    \bottomrule
   \end{tabular}
   \caption{#1}
   \label{tabla:#3}
  \end{center}
 \end{table}
}

%
% Nuevo comando para tablas pequeñas (menos de una página).
\newcommand{\tablaSmallSinColores}[5]{%
 \begin{table}[H]
  \begin{center}
   \begin{tabular}{#2}
    \toprule
    #4
    \otoprule
    #5
    \bottomrule
   \end{tabular}
   \caption{#1}
   \label{tabla:#3}
  \end{center}
 \end{table}
}

\newcommand{\tablaApaisadaSmall}[5]{%
\begin{landscape}
  \begin{table}
   \begin{center}
    \rowcolors {2}{gray!35}{}
    \begin{tabular}{#2}
     \toprule
     #4
     \otoprule
     #5
     \bottomrule
    \end{tabular}
    \caption{#1}
    \label{tabla:#3}
   \end{center}
  \end{table}
\end{landscape}
}

%
% Nuevo comando para tablas grandes con cabecera y filas alternas coloreadas en gris.
\newcommand{\tabla}[6]{%
  \begin{center}
    \tablefirsthead{
      \toprule
      #5
      \otoprule
    }
    \tablehead{
      \multicolumn{#3}{l}{\small\sl continúa desde la página anterior}\\
      \toprule
      #5
      \otoprule
    }
    \tabletail{
      \hline
      \multicolumn{#3}{r}{\small\sl continúa en la página siguiente}\\
    }
    \tablelasttail{
      \hline
    }
    \bottomcaption{#1}
    \rowcolors {2}{gray!35}{}
    \begin{xtabular}{#2}
      #6
      \bottomrule
    \end{xtabular}
    \label{tabla:#4}
  \end{center}
}

%
% Nuevo comando para tablas grandes con cabecera.
\newcommand{\tablaSinColores}[6]{%
  \begin{center}
    \tablefirsthead{
      \toprule
      #5
      \otoprule
    }
    \tablehead{
      \multicolumn{#3}{l}{\small\sl continúa desde la página anterior}\\
      \toprule
      #5
      \otoprule
    }
    \tabletail{
      \hline
      \multicolumn{#3}{r}{\small\sl continúa en la página siguiente}\\
    }
    \tablelasttail{
      \hline
    }
    \bottomcaption{#1}
    \begin{xtabular}{#2}
      #6
      \bottomrule
    \end{xtabular}
    \label{tabla:#4}
  \end{center}
}

%
% Nuevo comando para tablas grandes sin cabecera.
\newcommand{\tablaSinCabecera}[5]{%
  \begin{center}
    \tablefirsthead{
      \toprule
    }
    \tablehead{
      \multicolumn{#3}{l}{\small\sl continúa desde la página anterior}\\
      \hline
    }
    \tabletail{
      \hline
      \multicolumn{#3}{r}{\small\sl continúa en la página siguiente}\\
    }
    \tablelasttail{
      \hline
    }
    \bottomcaption{#1}
  \begin{xtabular}{#2}
    #5
   \bottomrule
  \end{xtabular}
  \label{tabla:#4}
  \end{center}
}



\definecolor{cgoLight}{HTML}{EEEEEE}
\definecolor{cgoExtralight}{HTML}{FFFFFF}

%
% Nuevo comando para tablas grandes sin cabecera.
\newcommand{\tablaSinCabeceraConBandas}[5]{%
  \begin{center}
    \tablefirsthead{
      \toprule
    }
    \tablehead{
      \multicolumn{#3}{l}{\small\sl continúa desde la página anterior}\\
      \hline
    }
    \tabletail{
      \hline
      \multicolumn{#3}{r}{\small\sl continúa en la página siguiente}\\
    }
    \tablelasttail{
      \hline
    }
    \bottomcaption{#1}
    \rowcolors[]{1}{cgoExtralight}{cgoLight}

  \begin{xtabular}{#2}
    #5
   \bottomrule
  \end{xtabular}
  \label{tabla:#4}
  \end{center}
}



\graphicspath{ {./img/} }

% Capítulos
\chapterstyle{bianchi}
\newcommand{\capitulo}[2]{
	\setcounter{chapter}{#1}
	\setcounter{section}{0}
	\setcounter{figure}{0}
	\setcounter{table}{0}
	\chapter*{#2}
	\addcontentsline{toc}{chapter}{#2}
	\markboth{#2}{#2}
}

% Apéndices
\renewcommand{\appendixname}{Apéndice}
\renewcommand*\cftappendixname{\appendixname}

\newcommand{\apendice}[1]{
	%\renewcommand{\thechapter}{A}
	\chapter{#1}
}

\renewcommand*\cftappendixname{\appendixname\ }

% Formato de portada
\makeatletter
\usepackage{xcolor}
\newcommand{\tutor}[1]{\def\@tutor{#1}}
\newcommand{\course}[1]{\def\@course{#1}}
\definecolor{cpardoBox}{HTML}{E6E6FF}
\def\maketitle{
  \null
  \thispagestyle{empty}
  % Cabecera ----------------
\noindent\includegraphics[width=\textwidth]{cabecera}\vspace{1cm}%
  \vfill
  % Título proyecto y escudo informática ----------------
  \colorbox{cpardoBox}{%
    \begin{minipage}{.8\textwidth}
      \vspace{.5cm}\Large
      \begin{center}
      \textbf{TFG del Grado en Ingeniería Informática}\vspace{.6cm}\\
      \textbf{\LARGE\@title{}}
      \end{center}
      \vspace{.2cm}
    \end{minipage}

  }%
  \hfill\begin{minipage}{.20\textwidth}
    \includegraphics[width=\textwidth]{escudoInfor}
  \end{minipage}
  \vfill
  % Datos de alumno, curso y tutores ------------------
  \begin{center}%
  {%
    \noindent\LARGE
    Presentado por \@author{}\\ 
    en Universidad de Burgos --- \@date{}\\
    Tutor: \@tutor{}\\
  }%
  \end{center}%
  \null
  \cleardoublepage
  }
\makeatother

\newcommand{\nombre}{Martín Encabo Contreras} %%% cambio de comando

% Datos de portada
\title{Creación de demos técnica y documentación asociada, para la detección/clasificación de materiales, utilizando un sensor radar a 60GHz}
\author{\nombre}
\tutor{nombre tutor}
\date{\today}

\begin{document}

\maketitle


\newpage\null\thispagestyle{empty}\newpage


%%%%%%%%%%%%%%%%%%%%%%%%%%%%%%%%%%%%%%%%%%%%%%%%%%%%%%%%%%%%%%%%%%%%%%%%%%%%%%%%%%%%%%%%
\thispagestyle{empty}


\noindent\includegraphics[width=\textwidth]{cabecera}\vspace{1cm}

\noindent D. nombre tutor, profesor del departamento de nombre departamento, área de nombre área.

\noindent Expone:

\noindent Que el alumno D. \nombre, con DNI 72897369L, ha realizado el Trabajo final de Grado en Ingeniería Informática titulado título de TFG. 

\noindent Y que dicho trabajo ha sido realizado por el alumno bajo la dirección del que suscribe, en virtud de lo cual se autoriza su presentación y defensa.

\begin{center} %\large
En Burgos, {\large \today}
\end{center}

\vfill\vfill\vfill

% Author and supervisor
\begin{minipage}{0.45\textwidth}
\begin{flushleft} %\large
Vº. Bº. del Tutor:\\[2cm]
D. nombre tutor
\end{flushleft}
\end{minipage}
\hfill
\begin{minipage}{0.45\textwidth}
\begin{flushleft} %\large
Vº. Bº. del co-tutor:\\[2cm]
D. nombre co-tutor
\end{flushleft}
\end{minipage}
\hfill

\vfill

% para casos con solo un tutor comentar lo anterior
% y descomentar lo siguiente
%Vº. Bº. del Tutor:\\[2cm]
%D. nombre tutor


\newpage\null\thispagestyle{empty}\newpage




\frontmatter

% Abstract en castellano
\renewcommand*\abstractname{Resumen}
\begin{abstract}
En este proyecto se va a desarrollar una aplicación mediante código Python que permita, mediante la lectura de materiales de un radar de 60 GHz, clasificar e indicar que tipo de material es. Además, se indica con que probabilidad es el material en cuestión.

Para realizar la parte del reconocimiento de materiales se aplicarán técnicas de aprendizaje automático, en concreto el clasificador Random Forest junto con la validación cruzada.

A la hora de recopilar la información necesaria, se usará la librería proporcionada por el fabricante del radar para conseguir una serie de lecturas con información suficiente capaz de identificar a cada objeto. Esta parte se realizará creando un registro de treinta materiales y diez lecturas de cada uno.

El objetivo del proyecto es el de demostrar el funcionamiento del sensor radar mediante una herramienta capaz de comparar diferentes tipos de materiales.

\end{abstract}

\renewcommand*\abstractname{Descriptores}
\begin{abstract}
Aprendizaje Automático, Clasificador, Random Forest, Reconocimiento, Python, Sensor, Radar, Acconeer
\end{abstract}

\clearpage

% Abstract en inglés
\renewcommand*\abstractname{Abstract}
\begin{abstract}
A \textbf{brief} presentation of the topic addressed in the project.
\end{abstract}

\renewcommand*\abstractname{Keywords}
\begin{abstract}
keywords separated by commas.
\end{abstract}

\clearpage

% Indices
\tableofcontents

\clearpage

\listoffigures

\clearpage

\listoftables
\clearpage

\mainmatter
\capitulo{1}{Introducción}

Descripción del contenido del trabajo y del estrucutra de la memoria y del resto de materiales entregados.

\capitulo{2}{Objetivos del proyecto}

En este apartado se explican los objetivos marcados en el proyecto, diferenciando entre los objetivos generales para llevar a cabo en este proyecto y los objetivos técnicos a lo largo de su ejecución.

\section{Objetivos generales}

A continuación se muestran los principales objetivos generales del proyecto:
\begin{itemize}
\item[•] El objetivo principal del proyecto se centra en documentar y demostrar el uso de un radar de 60GHz fabricado por Acconeer capaz de identificar, mediante un clasificador, el tipo de material al que pertenece un objeto.
\item[•] Crear un procedimiento de extracción de características y aprendizaje automático.

\item[•]Otro objetivo es el de mostrar un gráfico de probabilidades de la clasificación del material.

\item[•]Se intentará desarrollar una interfaz amigable para Windows.
\end{itemize}

\section{Objetivos técnicos}
\capitulo{3}{Conceptos teóricos}

Para construir un modelo capaz de capturar con precisión las características distintivas de las superficies objetivo, primero se debe comprender el origen y la estructura de la señal recibida. En este apartado se introducen algunos conceptos fundamentales del sistema de radar.
\section{Teledetección}

Teledetección o detección a distancia (detección remota) es la técnica que permite recopilar información a distancia de objetos sin que exista un contacto material. Para que esto sea posible tendremos una interacción entre los objetos y un sensor situado en una plataforma.

La distancia a que debe estar situado un sensor para ser remoto puede variar desde pocos decímetros hasta miles de kilómetros.

La teledetección es un flujo de radiación emitido por los objetos o materiales hacia un radar o sensor. El origen del flujo puede venir de:

\begin{itemize}
\item Radiación solar reflejada por los objetos - Teledetección pasiva
\item Radiación terrestre emitida por los objetos - Teledetección pasiva
\item Radiación emitida por el sensor y reflejada por los objetos (radar) - Teledetección activa
\end{itemize}

Se llama teledetección pasiva a la técnica aplicada por los sensores que miden las variaciones de la energía procedente de los objetos sin intervenir en el campo natural y se denomina teledetección activa a aquellos que generan un campo de energía artificial, utilizado para registrar y medir el efecto que en él producen los objetos.
 
Las características técnicas del sensor influye en la calidad de los datos y en la posibilidad de recibir información en distintas longitudes de onda. 

\imagen{sitemas_globales}{Sistemas de observación global.}

En la figura 3.1 vemos la globalización actual de la teledetección utilizada para obtener información precisa, actualizada y de fácil acceso empleada para diversas utilidades y finalidades como pueden ser:
\begin{itemize}
\item Imágenes territoriales
\item Motorización de las mareas
\item Análisis de materiales terrestres
\end{itemize}

\section{Radar Acconeer}

El radar utilizado en el proyecto está fabricado por Acconeer llamado A111. Es un radar de 60GHz basado en impulsos tecnología de radar coherente (PCR\footnote{Pulsed
Coherent Radar}) totalmente integrado en un pequeño chip de 29 mm2.
Esto permitirá una fácil integración en cualquier dispositivo portátil impulsado por batería.

\imagen{radar}{Radar A111.}

Aplicaciones:
\begin{itemize}
	\item Mediciones de distancia de alta precisión con mm de precisión y alta tasa de actualización.
	\item Detección de movimiento.
	\item Detección de velocidad.
	\item Detección de material.
	\item Seguimiento de objetos de alta precisión como el control de gestos.
	\item Seguimiento de alta precisión de objetos 3D.
	\item Control de los signos vitales, como la respiración y la frecuencia del pulso.
\end{itemize}

\imagen{figBlockDiagram}{Diagrama de bloques del sensor A111.}

La figura 3.2 muestra un diagrama de bloques del sensor A111. La señal se transmite desde la antena Tx y es recibida por la antena Rx, ambas integradas en la capa superior del sustrato del paquete A111. Además de la radio mmWave, el sensor consta de administración de energía y control digital, cuantificación de señales, memoria y un circuito de temporización.

El sensor se puede ejecutar en uno de los siguientes servicios básicos de la tabla 3.1.

\tablaSmall{Servicios del radar A111}{l c c}{serviciosa111}
{\multicolumn{1}{l}{Servicio} & Tipo de dato & Ejemplo de uso \\}{ 
Envelope & Amplitud & Distancia absoluta y
presencia estática \\
IQ & Amplitud y fase & Detección de obstáculos, respiración y distancia relativa \\
Sparse & Amplitud instantánea & Velocidad, detección de presencia y detección de gestos \\
} 


\section{Servicio IQ}

El servicio IQ utiliza la coherencia de fase del radar pulsado Acconeer para producir componentes estables en fase y en cuadratura. Este servicio se puede utilizar para la detección de presencia frente al sensor, la detección de la frecuencia respiratoria, la detección de obstáculos y , en nuestro caso, para diferenciar materiales.

Los componentes en fase y en cuadratura se representan como valores complejos, lo que genera un conjunto complejo de N\textsubscript{D} muestras representadas como x[d], dónde d es el índice de demora de la muestra.

Los datos obtenidos a través del servicio IQ proporcionan un método para examinar la reflectividad a diferentes distancias del sensor de radar.


\section{Aprendizaje automático}

El aprendizaje automático o machine learning es un tipo de inteligencia artificial (AI),
consiste en programar una computadora para que mejore en la realización de una tarea a partir de datos de ejemplo o de la experiencia.

\capitulo{4}{Técnicas y herramientas}

\section{Técnicas y metodologías}

\subsection{Metodología SCRUM}
Se trata de una metodología de trabajo ágil que tiene como finalidad dividir en periodos de tiempo el flujo de trabajo. Estos periodos son conocidos como  \textit{sprints}, al finalizar cada \textit{sprint} se realizan revisiones y reuniones donde se deciden las tareas de los próximos \textit{sprints}.

\subsection{Cliente de control de versiones}
\begin{itemize}
\item Herramientas consideradas: GitHub Desktop y Gitkraken
\item Herramienta elegida: GitHub
\end{itemize}

GitHub basado en Git junto con la extensión ZenHub se han utilizado para la gestión ágil del proyecto.

\subsection{Hosting del repositorio}
\begin{itemize}
\item Herramientas consideradas: GitLab y GitHub.
\item Herramienta elegida: GitHub
\end{itemize}

GitHub ofrece una gran cantidad de facilidades para mantener el proyecto en la nube y debido a que no cobra por sus servicios lo convierte en la mejor opción posible.


\section{Lenguajes y bibliotecas}

\subsection{Python}
Python es un lenguaje de programación que destaca por su código limpio y legible esto hace que sea uno de los lenguajes de iniciación de muchos programadores.
Además se trata de un lenguaje de multiparadigma y multiplataforma muy utilizado en la técnica del BigData.

\subsection{TensorFlow}
TensorFlow es una biblioteca de software de código abierto para computación numérica. Desarrollada por Google capaz de construir y entrenar redes neuronales.

\subsection{NumPy}
NumPy es una biblioteca utilizada en la programacion con Python para crear vectores y matrices grandes multidimensionales.

\subsection{Pandas}
Pandas es una biblioteca especializada en el manejo y análisis de estructuras de datos.

\subsection{acconeer.exptool}
acconeer.exptool es la biblioteca facilitada por la empresa Acconeer. Con ella haremos funcionar el radar de 60GHz.

\section{Herramientas de desarrollo}
\subsection{Jupyter Notebook}
Jupyter Notebook es el entorno de trabajo utilizado en el proyecto que permite desarrollar código en Python de manera dinámica. Nos ofrece integrar en un mismo archivo bloques de código, texto, gráficas o imágenes. Utilizado ampliamente en análisis numéricos y estadísticos.

\section{Herramientas de documentación}
\subsection{\LaTeX}
\LaTeX es un sistema de elaboración de documentos basado en comandos. El proyecto está desarrollado mediante una plantilla creada en \LaTeX.

\subsection{MiKTeX}
MiKTeX es una distribución de \LaTeX encargada de gestionar los componentes y paquetes. Tiene la capacidad de actualizarse así mismo descargando nuevas versiones de componentes.

\subsection{Texmaker}
Texmaker es el editor de documentos \LaTeX utilizado para funcionar necesita MiKTeX.

\capitulo{5}{Aspectos relevantes del desarrollo del proyecto}

Este apartado pretende recoger los aspectos más interesantes del desarrollo del proyecto, comentados por los autores del mismo.
Debe incluir desde la exposición del ciclo de vida utilizado, hasta los detalles de mayor relevancia de las fases de análisis, diseño e implementación.
Se busca que no sea una mera operación de copiar y pegar diagramas y extractos del código fuente, sino que realmente se justifiquen los caminos de solución que se han tomado, especialmente aquellos que no sean triviales.
Puede ser el lugar más adecuado para documentar los aspectos más interesantes del diseño y de la implementación, con un mayor hincapié en aspectos tales como el tipo de arquitectura elegido, los índices de las tablas de la base de datos, normalización y desnormalización, distribución en ficheros3, reglas de negocio dentro de las bases de datos (EDVHV GH GDWRV DFWLYDV), aspectos de desarrollo relacionados con el WWW...
Este apartado, debe convertirse en el resumen de la experiencia práctica del proyecto, y por sí mismo justifica que la memoria se convierta en un documento útil, fuente de referencia para los autores, los tutores y futuros alumnos.

\capitulo{6}{Trabajos relacionados}

En este apartado se muestran las herramientas relacionadas con el proyecto.
Hoy en día hay poca información y documentación de este tipo de tecnología de radares en la red, por ello solo muestro una única herramienta.

\section{\textit{RadarCat}}
\textit{RadarCat} (\textit{Radar Categorization for Input and Interaction}) permite que un dispositivo electrónico pueda reconocer y clasificar distintos materiales y objetos a tiempo real con alta precisión. Además da mucha más información de los objetos, como su estructura interna.

El escaner/radar utilizado envía ondas electromagnéticas al objeto, estas rebotan y vuelven al punto de partida para su procesamiento con aprendizaje automático o \textit{machine learning}.

Es necesario un ordenador para inicializar el software de \textit{RadarCat}, a este equipo se conecta una placa base donde se encuentran las antenas/radares capaces de realizar lecturas de objetos.
Por último encontramos una placa que protege el sensor y soporta el objeto a reconocer.

Más allá de la interacción humana con la computadora, \textit{RadarCat} también abre nuevas oportunidades en áreas como la navegación y el conocimiento mundial (por ejemplo, usuarios con baja visión).

Aquí se puede acceder a un vídeo interesante que muestra los múltiples usos de \textit{RadarCat} \url{https://www.youtube.com/watch?v=B6sn2vRJXJ4}

En el proyecto \textit{RadarCat} creado por \textit{Google} utiliza un radar del fabricante \textit{Infineon Technologies} al que \textit{Google} llama \textit{Google ATAP Project Soli}, nosotros para este proyecto estamos utilizando el del fabricante \textit{Acconeer}. Existe una gran diferencia de precios entre los dos, mientras que nuestro radar tiene un precio de 52,73€ el de \textit{Infineon}, según muestra la web del fabricante, llega a valer 284,64€.


\section{Exploración de interacciones tangibles con sensores de radar}

En esta sección se expone un artículo que muestra un ejemplo de uso para sensores de radar y aprendizaje automático para la exploración de interacciones tangibles\footnote{El mundo digital no se observa mediante un monitor o la pantalla de un dispositivo móvil, sino que está mezclado con la realidad.}.\cite{Yeo2018Dec}  

El el artículo se proponen usar el radar como una plataforma de detección para rastrear la identidad de un objeto o la cantidad de objetos próximos para habilitar la \textit{Tangible User Interface} (TUI). El radar mediante redes microondas ha sido capaz de reconocer, contar y ordenar objetos y además determinar movimientos (deslizamiento o rotación). Todo ello con elementos comunes, generalmente objetos en gran parte planos.

Para realizar el prototipo, se usa el sensor de radar \textit{Google ATAP Project Soli}, por lo que este proyecto es una ampliación de \textit{RadarCat}. Para mejorar el objetivo buscado se crean unas estructuras 3D, mediante una impresora, para guiar la onda del radar y mejorar la colocación de los objetos sobre este. Según indica el articulo se han creado hasta 12 estructuras distintas (p.ej. soporte para fichas \textit{Lego}, tarjetero y un espacio de discos), cada una para un fin.

El prototipo se basa en el paradigma \textit{Token+Constraint}. En estas interfaces, los \textit{tokens} (fichas) son objetos físicos discretos que representan información digital. Las \textit{constraint} (restricciones) son regiones de confinamiento que se asignan a las operaciones digitales. Estos se incorporan con frecuencia como estructuras que canalizan mecánicamente cómo se pueden manipular las fichas, a menudo limitando su movimiento a un solo grado de libertad. La colocación y manipulación de tokens dentro de sistemas de restricciones se puede usar para invocar y controlar una variedad de interpretaciones computacionales. Se muestra que al agregar un «caso guía» que
actua como restriccion, se puede ampliar en gran medida las capacidades de interacción y mejorar la precisión de detección de objetos.

En el siguiente enlace se puede acceder a un vídeo donde se muestran varios ejemplos de uso del prototipo: \url{https://www.youtube.com/watch?v=zDAJ_0KIovc}


\capitulo{7}{Conclusiones y Líneas de trabajo futuras}

Todo proyecto debe incluir las conclusiones que se derivan de su desarrollo. Éstas pueden ser de diferente índole, dependiendo de la tipología del proyecto, pero normalmente van a estar presentes un conjunto de conclusiones relacionadas con los resultados del proyecto y un conjunto de conclusiones técnicas. 
Además, resulta muy útil realizar un informe crítico indicando cómo se puede mejorar el proyecto, o cómo se puede continuar trabajando en la línea del proyecto realizado. 



\bibliographystyle{ieeetr}
\bibliography{bibliografia}

\end{document}
